\documentclass[11pt,a4paper]{article}
\usepackage[margin=1.8cm,headheight=48pt,headsep=12pt,footskip=22pt]{geometry}
\usepackage{fontspec}
\setmainfont{Noto Serif}
\setsansfont{Noto Sans}
\setmonofont{DejaVu Sans Mono}
\usepackage{amsmath,amssymb,mathtools}
\usepackage{enumitem,tabularx,booktabs,array,multicol}
\usepackage{xcolor,graphicx,tikz,fancyhdr,lastpage}
\usepackage{listings,tcolorbox,hyperref}
\hypersetup{colorlinks=true,linkcolor=blue!60!black,urlcolor=blue!60!black}
\definecolor{hcmusblue}{RGB}{28,55,111}
\definecolor{lightblue}{RGB}{237,242,250}
\definecolor{codebg}{RGB}{247,248,250}
\setlength{\parindent}{0pt}
\setlength{\parskip}{4pt}
\setlist[enumerate]{leftmargin=*,itemsep=5pt,topsep=4pt}
\setlist[itemize]{leftmargin=*,itemsep=2pt,topsep=2pt}
\lstset{language=C++,basicstyle=\ttfamily\footnotesize,backgroundcolor=\color{codebg},frame=single,rulecolor=\color{black!15},breaklines=true,showstringspaces=false,columns=fullflexible,keywordstyle=\color{hcmusblue}\bfseries,commentstyle=\color{green!40!black},numbers=left,numberstyle=\tiny\color{black!45},xleftmargin=1.5em,framexleftmargin=1.2em}
\pagestyle{fancy}
\fancyhf{}
\fancyhead[L]{\includegraphics[height=1.05cm,keepaspectratio]{../../04_Nhan_dien/hcmus_logo.png}\hspace{0.45em}\begin{minipage}[c]{0.64\textwidth}\textbf{TRƯỜNG ĐẠI HỌC KHOA HỌC TỰ NHIÊN}\\[-1pt]\small ĐHQG TP. HỒ CHÍ MINH\end{minipage}}
\fancyhead[R]{\small\textcolor{hcmusblue}{TÀI LIỆU LUYỆN TẬP}}
\fancyfoot[L]{\scriptsize Biên soạn độc lập - Không phải đề thi chính thức}
\fancyfoot[R]{\scriptsize Trang \thepage}
\renewcommand{\headrulewidth}{0.5pt}
\renewcommand{\footrulewidth}{0.3pt}
\newcommand{\examtitle}[4]{\begin{center}{\Large\bfseries\color{hcmusblue}#1}\par\vspace{3pt}{\large\bfseries #2}\par\vspace{3pt}\small #3\quad\textbullet\quad #4\end{center}\vspace{5pt}\hrule\vspace{8pt}}
\newcommand{\studentline}{\begin{tabularx}{\textwidth}{X X}\textbf{Họ và tên:} \hrulefill & \textbf{MSSV:} \hrulefill\\[7pt]\textbf{Lớp:} \hrulefill & \textbf{Ngày làm:} \hrulefill\end{tabularx}\vspace{8pt}}
\newcommand{\points}[1]{\hfill\textbf{(#1 điểm)}}
\newtcolorbox{notebox}{colback=lightblue,colframe=hcmusblue!55,boxrule=0.5pt,arc=1.5mm,left=2mm,right=2mm,top=1mm,bottom=1mm}
\begin{document}
\small
\examtitle{ĐÁP ÁN ĐỀ LUYỆN TẬP 04 - CƠ SỞ LẬP TRÌNH}{CƠ SỞ LẬP TRÌNH - HCMUS (theo lịch học tuần 1--5)}{Lời giải chi tiết và mã C++11}{Giữa kỳ 1 - HK1 2026--2027}

\section*{Đáp án và lời giải - Đề 04}
\begin{enumerate}
\item \begin{minipage}[t]{\linewidth}
Xét điều kiện từ trên xuống dưới để tránh nhánh rộng hơn bắt mất trường hợp cụ thể:
\begin{lstlisting}
if (x < 0 || x > 10) cout << "Invalid";
else if (x >= 8.5) cout << "Excellent";
else if (x >= 7.0) cout << "Good";
else if (x >= 5.0) cout << "Pass";
else cout << "Not pass";
\end{lstlisting}
\end{minipage}
\item \begin{minipage}[t]{\linewidth}
Mã C++11 tham khảo:
\begin{lstlisting}
int countDivisors(int n) {
    int count = 0;
    for (int d = 1; d <= n; ++d)
        if (n % d == 0) ++count;
    return count;
}
\end{lstlisting}
Ví dụ $n=12$ có các ước $1,2,3,4,6,12$, kết quả 6.
\end{minipage}
\item \begin{minipage}[t]{\linewidth}
Mã C++11 tham khảo:
\begin{lstlisting}
const double PI = 3.141592653589793;
double area(double r) {
    return PI * r * r;
}
double area(double w, double h) {
    return w * h;
}
\end{lstlisting}
Khác số lượng tham số nên trình biên dịch phân biệt được hai hàm cùng tên.
\end{minipage}
\item \begin{minipage}[t]{\linewidth}
Mã C++11 tham khảo:
\begin{lstlisting}
void moveZerosToEnd(int a[], int n) {
    int write = 0;
    for (int read = 0; read < n; ++read)
        if (a[read] != 0) a[write++] = a[read];
    while (write < n) a[write++] = 0;
}
\end{lstlisting}
Ví dụ $[0,4,0,2,7]$ trở thành $[4,2,7,0,0]$. Khi ghi, chỉ ghi vào vị trí không vượt quá chỉ số đang đọc nên thứ tự phần tử khác 0 được bảo toàn.
\end{minipage}
\item \begin{minipage}[t]{\linewidth}
Mã C++11 tham khảo:
\begin{lstlisting}
long long borderSum(const int a[][MAX], int R, int C) {
    long long sum = 0;
    for (int i = 0; i < R; ++i)
        for (int j = 0; j < C; ++j)
            if (i == 0 || i == R-1 || j == 0 || j == C-1)
                sum += a[i][j];
    return sum;
}
\end{lstlisting}
Mỗi ô được duyệt đúng một lần và kiểm tra điều kiện biên bằng phép OR, do đó không đếm trùng các góc hoặc các ô khi $R=1$ hay $C=1$.
\end{minipage}
\end{enumerate}

\end{document}
